\section{Classical recovery from evidential regularity}
\label{sec:classical-recovery}

The preceding matroid analysis raises a natural reverse question.  Four-valued
semantics permits both evidential gaps and evidential gluts; must classical logic
therefore be rejected globally, or can a classical fragment be recovered when the
underlying evidence structure satisfies stronger regularity conditions?  This section
formalizes the latter possibility at the coarse two-channel level.

The result is deliberately local and structural.  I do not assume bivalence as an
extra semantic postulate.  Instead, the positive and negative evidence channels are
constrained, and the induced four-valued state is then classified.  The resulting
recovery separates three ideas that are easily conflated: local classical value,
recovery of individual classical inference principles, and persistence of classicality
under modal or evidential change.

\subsection{Completeness and consistency of the channel geometry}

Let \(A\) be a body of evidence with polarity map
\(\chi:E\to\{+,-\}\).  Write
\[
  S_+(A) \quad\text{and}\quad S_-(A)
\]
for support of the positive and negative channels.  Define
\begin{align*}
  \mathsf{Complete}(A)
    &:\Longleftrightarrow S_+(A)\lor S_-(A),\\
  \mathsf{Consistent}(A)
    &:\Longleftrightarrow \neg\bigl(S_+(A)\land S_-(A)\bigr),\\
  \mathsf{Regular}(A)
    &:\Longleftrightarrow
       \mathsf{Complete}(A)\land\mathsf{Consistent}(A).
\end{align*}
At the value level let \(\mathsf{GapFree}(v)\) mean that \(v\neq\Nval\) in the
bilateral sense, and \(\mathsf{GlutFree}(v)\) that \(v\neq\Bval\) in the bilateral
sense.  The Lean definitions use the already existing Boolean gap/glut tests rather
than introducing a second four-valued algebra.

\begin{theorem}[Geometry--value recovery]
\label{thm:classical-recovery-geometry}
Suppose \(A\) realizes the FDE value \(v\).  Then
\begin{align*}
  \mathsf{GapFree}(v)
    &\Longleftrightarrow \mathsf{Complete}(A),\\
  \mathsf{GlutFree}(v)
    &\Longleftrightarrow \mathsf{Consistent}(A),\\
  \mathsf{Regular}(A)
    &\Longleftrightarrow v\in\{\Tval,\Fval\}.
\end{align*}
\statusverified
\end{theorem}

Thus the classical two-valued slice need not be imposed by deleting \(\Nval\) and
\(\Bval\) from the semantics.  It is exactly the image of evidence bodies satisfying
both coverage and non-conflict at the coarse channel level:
\[
  \boxed{
  \mathsf{Complete}+\mathsf{Consistent}
  \quad\Longleftrightarrow\quad
  \{\Tval,\Fval\}.}
\]
This is a recovery theorem inside the existing four-valued semantics, not a claim
that all evidence bodies satisfy its hypotheses.

\subsection{The recovered propositional slice}

The recovered values are closed under the existing FDE propositional operations.

\begin{proposition}[Classical subalgebra at the value level]
\label{prop:classical-recovery-subalgebra}
If \(v,w\in\{\Tval,\Fval\}\), then
\[
  \neg v\in\{\Tval,\Fval\},\qquad
  v\land w\in\{\Tval,\Fval\},\qquad
  v\lor w\in\{\Tval,\Fval\}.
\]
\statusverified
\end{proposition}

The proposition is modest but important.  Recovery is not merely a pointwise
classification of two values: the existing four-valued connectives restrict to a
closed classical propositional fragment on that slice.

\subsection{LEM and EFQ return for different reasons}

The two most conspicuous departures from classical reasoning separate cleanly under
the bilateral semantics.  Let
\[
  \operatorname{LEM}(v)=v\lor\neg v,
  \qquad
  \operatorname{Contr}(v)=v\land\neg v.
\]
For LP-style designation, a value is tolerated when its positive bit is present.

\begin{proposition}[Separated recovery conditions]
\label{prop:lem-efq-separated-recovery}
For every FDE value \(v\):
\begin{enumerate}[label=(\roman*)]
  \item \(\operatorname{LEM}(v)\) is LP-designated exactly when \(v\) is gap-free;
  \item \(\operatorname{LEM}(v)=\Tval\) exactly when
        \(v\in\{\Tval,\Fval\}\);
  \item if \(v\) is glut-free, then
        \(\operatorname{Contr}(v)\) LP-entails every conclusion value.
\end{enumerate}
\statusverified
\end{proposition}

Consequently, excluded middle and explosion do not reappear as an inseparable
``classical package.''  Gap-freedom is the relevant structural boundary for tolerant
excluded middle, whereas glut-freedom is the relevant boundary for LP ex falso.
Full local classicality requires both.  In schematic form,
\[
\begin{array}{ccl}
  \text{no gap} &\Longrightarrow& \text{tolerant LEM},\\
  \text{no glut} &\Longrightarrow& \text{LP EFQ},\\
  \text{no gap + no glut} &\Longleftrightarrow& \Tval/\Fval\text{ slice}.
\end{array}
\]
The LP-explosion statement is vacuous in the precise semantic sense that, under
GlutFree, the contradiction premise is not designated.  It should therefore not be
read as a proof that an inconsistent evidential state has somehow generated arbitrary
content.

\subsection{Stability is persistence, not local bivalence}

The modal-probabilistic refinement already distinguishes local status from stability
through an accessibility range.  This distinction survives the recovery analysis.
Local completeness and consistency are sufficient for a local \(\Tval\) or \(\Fval\)
value; no modal stability assumption is needed for that classification.

\begin{proposition}[Stable propagation of recovered classicality]
\label{prop:stable-classical-recovery}
Let \(v:W\to\{\Nval,\Tval,\Fval,\Bval\}\).  If \(v\) is stable at \(w\) and
\(v(w)\) is evidentially regular, then every world accessible from \(w\) carries a
classical value.  In fact stability gives the same classical value throughout the
accessible range.
\statusverified
\end{proposition}

This motivates a three-step reading of the original hypothesis:
\[
  \boxed{
  \text{coverage}\;\to\;\text{no gap},\qquad
  \text{consistency}\;\to\;\text{no glut},\qquad
  \text{stability}\;\to\;\text{persistence}.}
\]
The third condition strengthens classicality dynamically or modally rather than
creating local bivalence in the first place.

\subsection{Deletion robustness and fragility}

Gate~4 makes it possible to ask whether recovered classicality survives structural
changes in its evidence base.  The answer already distinguishes redundancy from
fragility.

\begin{theorem}[Parallel deletion preserves regularity]
\label{thm:parallel-deletion-preserves-regularity}
Suppose \(x\neq y\) are evidence atoms of the same polarity and \(y\) remains in the
ground set.  Deleting \(x\) preserves channel completeness, channel consistency, and
therefore channel regularity exactly:
\[
  \mathsf{Regular}(E\setminus\{x\})
  \Longleftrightarrow
  \mathsf{Regular}(E).
\]
\statusverified
\end{theorem}

Thus local classical recovery is robust to deletion of a redundant parallel source.
By contrast, the last source carrying a classical channel is structurally critical.

\begin{theorem}[Unique-source deletion destroys completeness]
\label{thm:unique-deletion-classical-to-gap}
Suppose \(E\) is channel-regular and \(c\) is the unique representative of its
polarity.  Then
\[
  \neg\mathsf{Complete}(E\setminus\{c\}).
\]
Moreover, every FDE value realized by the remaining evidence is \(\Nval\).
\statusverified
\end{theorem}

This theorem gives a precise sense in which a classical state can be genuine but
fragile.  The state is classical because the current evidence geometry is complete
and consistent, yet removal of a structurally indispensable source reopens an
evidential gap:
\[
  \Tval\to\Nval
  \qquad\text{or}\qquad
  \Fval\to\Nval.
\]
Classicality is therefore not merely a label on a value; its robustness depends on
redundancy in the supporting evidence geometry.

\subsection{Independent-capacity recovery and contraction}

Gate~4 also showed that ordinary ground occupancy and residual independent-information
capacity diverge after contraction.  This motivates a stronger recovery predicate.
Call a channel \emph{nonloop-supported} when it has a remaining ground atom that is
not a loop in the current matroid minor.  Define \(\mathsf{NLComplete}\),
\(\mathsf{NLConsistent}\), and \(\mathsf{NLRegular}\) by replacing ordinary channel
support with nonloop support in the preceding definitions.

\begin{theorem}[Occupancy and capacity coincide before minors]
\label{thm:nonloop-regularity-before-minors}
For the original finite channel matroid,
\[
  \mathsf{NLRegular}(M)=\mathsf{Regular}(E).
\]
Equivalently, nonloop completeness and consistency coincide separately with ordinary
channel completeness and consistency.
\statusverified
\end{theorem}

This equivalence holds because the canonical finite channel matroid is loopless before
minors.  Hence the original \(\Tval/\Fval\) recovery theorem can be read either as a
statement about occupancy or about independent channel capacity.

Contraction breaks that identification.

\begin{theorem}[Contraction destroys independent classical capacity]
\label{thm:contraction-destroys-nonloop-regularity}
Let \(E\) be channel-regular and let \(c\in E\).  After contracting \(c\), the
resulting minor is not nonloop-complete and therefore not nonloop-regular:
\[
  \neg\mathsf{NLComplete}(M/c),
  \qquad
  \neg\mathsf{NLRegular}(M/c).
\]
\statusverified
\end{theorem}

The reason is structural.  Contraction makes every remaining atom of \(c\)'s polarity
a loop, while channel regularity already excluded support for the opposite polarity.
Thus no independent channel direction remains.

The distinction is observable even when a same-polarity atom survives in the ground
set.

\begin{proposition}[Occupancy without capacity]
\label{prop:occupancy-without-capacity}
Suppose \(c\neq y\) are ground atoms of the same polarity.  After contracting \(c\),
that polarity is still ordinarily supported by \(y\), but it has no nonloop support.
\statusverified
\end{proposition}

Therefore a coarse support reading can continue to display a channel that has become
structurally informationally inert.  Gate~5 consequently distinguishes
\[
  \boxed{
  \text{classical-valued occupancy}
  \quad\neq\quad
  \text{classicality sustained by residual independent capacity}.}
\]
This is the first formal refinement of ``informational stability'' in the recovery
program: classical value recovery is local, whereas structural classicality must also
track what independent information survives admissible transformations.

\subsection{Scope and next boundary}

The current Gate~5 results establish value-level, modal-persistence, deletion, and
independent-capacity recovery layers.  They do \emph{not} yet prove completeness of a
classical proof calculus, classicality of every modal formula, or preservation under
every evidence update.  In particular, the present LP ex-falso theorem is a
value-level semantic recovery under glut-freedom; its premise is undesignated there,
and no stronger proof-theoretic conclusion is inferred.

The reliability refinement raises the next boundary.  A coarse state can be regular
while distinct fine reliability profiles project to it.  A later gate should ask
whether classical recovery commutes with forgetting reliability and whether a
fine-grained notion of independent-capacity regularity survives that projection.
Likewise, a formula-level recovery theorem should quantify regularity over the atomic
valuation and prove closure inductively through the language, treating modal operators
separately from the already verified propositional \(\Tval/\Fval\) subalgebra.
